\chapter*{Appendix L: Spectral Methods}
\addcontentsline{toc}{chapter}{Appendix L: Spectral Methods}

\section{Spectral Decomposition of Operators}

\subsection{Self-Adjoint Operators}

Let $(\mathcal{H},\langle\cdot,\cdot\rangle)$ be a separable Hilbert space.  
A bounded operator $A:\mathcal{H}\to\mathcal{H}$ is self-adjoint if
\[
\langle Ax,y\rangle = \langle x,Ay\rangle
\quad
\text{for all } x,y\in\mathcal{H}.
\]

\begin{theorem}[Spectral Theorem]
If $A$ is self-adjoint, then there exists a projection-valued measure $E$ on $\mathbb{R}$ such that
\[
A = \int_{\mathbb{R}} \lambda \, dE(\lambda).
\]
\end{theorem}

\subsection{Functional Calculus}

For Borel function $f:\mathbb{R}\to\mathbb{C}$,
\[
f(A) = \int_{\mathbb{R}} f(\lambda) \, dE(\lambda).
\]

\section{Laplacian Spectra}

\subsection{Laplace--Beltrami Operator}

For a smooth manifold $(\mathcal{M},g)$, the Laplace--Beltrami operator is
\[
\Delta f = \nabla_i\nabla^i f.
\]

\begin{definition}[Spectrum]
The spectrum of $\Delta$ is the set of $\lambda\in\mathbb{R}$ for which
\[
\Delta f + \lambda f = 0
\]
has a nontrivial solution $f\in C^\infty(\mathcal{M})$.
\end{definition}

On compact $\mathcal{M}$, the spectrum is discrete:
\[
0 = \lambda_0 < \lambda_1 \le \lambda_2 \le \cdots,\quad \lambda_k\to\infty.
\]

\subsection{Eigenfunction Expansion}

Any square-integrable function $u\in L^2(\mathcal{M})$ admits the expansion
\[
u(x)
  = \sum_{k=0}^\infty \hat{u}_k \phi_k(x),
\qquad
\hat{u}_k = \int_{\mathcal{M}} u(x)\phi_k(x)\, dV_g.
\]

\section{Heat Kernel and Green's Functions}

\subsection{Heat Kernel}

The heat kernel $K_t(x,y)$ is defined as the fundamental solution to
\[
\partial_t u = \Delta u,
\quad
u(0,x) = f(x).
\]
It satisfies
\[
K_t(x,y)
 = \sum_{k=0}^\infty e^{-\lambda_k t} \phi_k(x)\phi_k(y).
\]

\subsection{Green's Function}

For the Poisson equation $-\Delta u = f$, a Green's function $G(x,y)$ satisfies
\[
-\Delta_x G(x,y) = \delta_y(x).
\]
Spectrally,
\[
G(x,y)
  = \sum_{k=1}^\infty \frac{1}{\lambda_k} \phi_k(x)\phi_k(y),
\]
excluding the constant eigenfunction when $\lambda_0=0$.

\section{Fourier and Wavelet Bases}

\subsection{Fourier Basis on the Torus}

On $\mathbb{T}^n = (\mathbb{R}/2\pi\mathbb{Z})^n$,
\[
e_k(x) = \frac{1}{(2\pi)^{n/2}} e^{i k\cdot x},
\qquad k\in\mathbb{Z}^n,
\]
form an orthonormal basis diagonalising $\Delta$:
\[
\Delta e_k = -|k|^2 e_k.
\]

\subsection{Wavelet Bases}

A wavelet basis is a collection $\{\psi_{j,k}\}$ with scale $j$ and shift $k$ such that
\[
f(x) = \sum_{j,k} c_{j,k} \psi_{j,k}(x)
\]
converges for all $f\in L^2(\mathbb{R}^n)$.
Wavelets enable multiscale spectral representations.

\section{Spectral Geometry}

\subsection{Spectral Invariants}

Key curvature-sensitive quantities include:

\begin{itemize}
\item Heat trace:
\[
Z(t) = \sum_{k=0}^\infty e^{-\lambda_k t}.
\]
\item Weyl asymptotics:
\[
N(\lambda)
  \sim \frac{\omega_n}{(2\pi)^n}
       \operatorname{Vol}(\mathcal{M})\, \lambda^{n/2}.
\]
\end{itemize}

\subsection{Spectral Distance}

For two manifolds $(\mathcal{M}_1,g_1)$ and $(\mathcal{M}_2,g_2)$, define
\[
d_{\mathrm{spec}}^2
  = \sum_{k=0}^\infty
    w_k \left(\lambda_k^{(1)} - \lambda_k^{(2)}\right)^2,
\]
with weights $w_k$ decaying sufficiently fast.

\section{Discrete Spectral Methods}

\subsection{Graph Laplacian}

For a weighted graph with adjacency $w_{ij}$ and degree $d_i$,
\[
L_{ij} = \begin{cases}
 d_i, & i=j,\\
 -w_{ij}, & i\ne j.
\end{cases}
\]

\subsection{Spectral Clustering}

Compute the eigenvectors associated with the $k$ smallest eigenvalues of $L$ and cluster their rows in $\mathbb{R}^k$.

\subsection{Finite Element Spectra}

Given finite element basis $\{\phi_i\}$ on mesh $\mathcal{T}_h$,
\[
K_{ij} = \int_{\mathcal{M}} \nabla\phi_i\cdot\nabla\phi_j\, dV_g,
\qquad
M_{ij} = \int_{\mathcal{M}} \phi_i \phi_j \, dV_g.
\]

Solve the generalised eigenvalue problem:
\[
K u = \lambda M u.
\]

\section{Spectral Filtering}

\subsection{Low-Pass and High-Pass Operators}

For filter function $\sigma(\lambda)$,
\[
\sigma(\Delta) f
  = \sum_{k=0}^\infty \sigma(\lambda_k) \hat{f}_k \phi_k.
\]

Common choices:
\begin{itemize}
\item Low-pass: $\sigma(\lambda)=e^{-\lambda/\Lambda}$  
\item High-pass: $\sigma(\lambda)=1-e^{-\lambda/\Lambda}$  
\item Band-pass: indicator functions of intervals  
\end{itemize}

\section{Pseudo-Differential Symbols}

\subsection{Symbol of an Operator}

A pseudo-differential operator $A$ on $\mathbb{R}^n$ has symbol $a(x,\xi)$ such that
\[
Af(x)
  = \frac{1}{(2\pi)^n}
    \int_{\mathbb{R}^n}
    e^{i x\cdot\xi} a(x,\xi)\widehat{f}(\xi)\, d\xi.
\]

\subsection{Ellipticity}

$A$ is elliptic of order $m$ if
\[
|a(x,\xi)| \ge C |\xi|^m
\quad
\text{for } |\xi|\gg 1.
\]

\section{Spectral Evolution Equations}

\subsection{Heat Equation}

The evolution in spectral form:
\[
u(t,x)
  = \sum_{k=0}^\infty e^{-\lambda_k t} \hat{u}_k(0) \phi_k(x).
\]

\subsection{Wave Equation}

\[
u(t,x)
  = \sum_{k=0}^\infty
    \left(
    a_k \cos(\sqrt{\lambda_k} t)
    + b_k \frac{\sin(\sqrt{\lambda_k} t)}{\sqrt{\lambda_k}}
    \right)
    \phi_k(x).
\]

\subsection{Diffusion Through Operators}

For any generator $A$ with spectrum $\{\mu_k\}$,
\[
e^{tA} f = \sum_k e^{t\mu_k} \hat{f}_k \phi_k.
\]


